\documentclass[11pt,a4paper]{article}
\usepackage[T1]{fontenc}
\usepackage[utf8]{inputenc}
\usepackage[czech]{babel}
\usepackage{lmodern,microtype}
\usepackage{amsmath,amssymb,mathtools}
\usepackage{geometry}
\geometry{left=19mm,right=19mm,top=21mm,bottom=21mm}
\usepackage[dvipsnames]{xcolor}
\usepackage{enumitem,tabularx,booktabs}
\usepackage[most]{tcolorbox}
\usepackage{fancyhdr,lastpage}
\usepackage[hidelinks]{hyperref}
\definecolor{deep}{HTML}{203955}
\definecolor{teal}{HTML}{176E76}
\definecolor{slate}{HTML}{506173}
\definecolor{pale}{HTML}{F3F7F9}
\setlength{\parindent}{0pt}
\setlength{\parskip}{5pt}
\setlist[enumerate]{leftmargin=6mm,itemsep=5pt,topsep=3pt}
\pagestyle{fancy}\fancyhf{}
\fancyhead[L]{\sffamily\footnotesize\color{slate} MAT · 4. ročník · kombinatorika}
\fancyhead[R]{\sffamily\footnotesize\color{slate} Hodina \LESSONNUM\, / 112}
\fancyfoot[L]{\sffamily\footnotesize\color{slate} Výukový podklad · samostatný dokument}
\fancyfoot[R]{\sffamily\footnotesize\color{slate} \thepage/\pageref{LastPage}}
\renewcommand{\headrulewidth}{0.35pt}
\renewcommand{\footrulewidth}{0pt}
\newcommand{\sect}[1]{\par\vspace{2mm}{\Large\sffamily\bfseries\color{deep} #1}\par\vspace{-0.5mm}\noindent\color{teal}\rule{\linewidth}{0.7pt}\color{black}\vspace{1mm}}
\tcbset{theobox/.style={enhanced,breakable,colback=pale,colframe=teal!45!white,boxrule=.6pt,arc=1mm,left=3mm,right=3mm,top=2mm,bottom=2mm}}
\newcommand{\LESSONNUM}{33}
\begin{document}

{\Huge\sffamily\bfseries\color{deep} Permutace s opakováním}\par
{\sffamily\large\color{teal} Výukový podklad · 33. vyučovací hodina}\par
\textbf{Hodinová dotace:} 45 minut\quad \textbf{Ročník:} 4. SPŠ\quad \textbf{Rozpis:} TP 2.1\par
\begin{tcolorbox}[theobox]\textbf{Výukový cíl.} Počítat pořadí multisetů se stejnými prvky, odvodit dělení faktoriály četností a vyřešit jednoduchá omezení.\end{tcolorbox}
\sect{Teorie a vysvětlení}
Pokud ve skupině \(n\) prvků rozlišujeme pouze typy, které se vyskytují \(n_1,n_2,\ldots,n_r\)-krát, kde \(n_1+\cdots+n_r=n\), počet různých pořadí je
\[P(n;n_1,\ldots,n_r)=\frac{n!}{n_1!\,n_2!\cdots n_r!}.\]
Důvod: při hypotetickém rozlišení všech kopií napočítáme \(n!\) pořadí, ale výměny \(n_i\) nerozlišitelných kopií žádné nové pořadí netvoří; každé skutečné pořadí jsme proto započetli \(n_1!\cdots n_r!\)-krát.\par
\textbf{Rozhodovací znak:} uspořádáme všechny prvky, ale některé jsou stejné. Nevkládej vzorec pro variace s opakováním: tam se naopak znaky mohou \emph{svobodně volit} z celé abecedy na každé pozici.
\sect{Řešené příklady}
\textbf{Příklad 1.} Kolik různých slov (nemusejí dávat smysl) vznikne přeskupením písmen \textsc{ANANAS}?\par
\textit{Řešení.} Je \(6\) písmen: \(A\) třikrát, \(N\) dvakrát, \(S\) jednou. Počet \(\frac{6!}{3!\,2!}=\frac{720}{12}=60\).\par\smallskip
\textbf{Příklad 2.} Kolik pořadí vznikne z písmen \textsc{KOKOS}?\par
\textit{Řešení.} Pět písmen; \(K\) dvakrát, \(O\) dvakrát, \(S\) jednou. \(\frac{5!}{2!\,2!}=30\).\par\smallskip
\textbf{Příklad 3.} Kolik sedmimístných řetězců obsahuje právě čtyři nuly a tři jedničky?\par
\textit{Řešení.} Vybíráme, kam postavíme tři jedničky; ostatní pozice budou nuly. \(\frac{7!}{4!\,3!}=35\).\par\smallskip
\textbf{Příklad 4.} U slova \textsc{ANANAS} požadujeme, aby všechna tři \(A\) stála vedle sebe.\par
\textit{Řešení.} Blok \(AAA\), dále \(N,N,S\): čtyři objekty se dvěma nerozlišitelnými \(N\). Počet \(\frac{4!}{2!}=12\). Blok \(AAA\) uvnitř nemá další různá pořadí.\par\smallskip
\textbf{Příklad 5.} Kolik různých rozložení pěti shodných červených a dvou shodných modrých žetonů do sedmi míst v řadě?\par
\textit{Řešení.} Sedm pozic, z toho dvě pro modré: \(\frac{7!}{5!\,2!}=21\).\par\smallskip
\clearpage\sect{Procvičování pro studenty}
\begin{enumerate}
\item Přeskup písmena \textsc{TATRA}.
\item Kolik různých řad tvoří tři nuly a dvě jedničky?
\item Přeskup \textsc{MAMA} tak, aby všechna \(M\) byla vedle sebe.
\item Kolik pořadí má osm žetonů: \(3\) červené, \(3\) modré a \(2\) zelené?
\item Ve slově \textsc{ANANAS} mají být obě \(N\) vedle sebe.
\end{enumerate}
\sect{Řešení a zdůvodnění k procvičování}
\begin{enumerate}
\item Pět znaků, \(T\) dvakrát, \(A\) dvakrát: \(5!/(2!2!)=30\).
\item \(5!/(3!2!)=10\).
\item Blok \(MM\) a dvě \(A\): \(3!/2!=3\) možnosti.
\item \(8!/(3!3!2!)=560\).
\item Blok \(NN\), tři \(A\) a \(S\): \(5!/3!=20\).
\end{enumerate}
\sect{Metodické poznámky}
\textbf{Typické chyby:}
\begin{itemize}[leftmargin=5mm,itemsep=1pt]
\item Dělíme faktoriály četností, nikoli samotnými četnostmi.
\item Nerozlišitelné prvky nemají vnitřní \(2!\) uspořádání.
\item Při podmínce „vedle sebe“ přepočítat počet nových objektů.
\end{itemize}
\textbf{Doporučený průběh 45 minut:} 0–10 motivace na slovech; 10–20 odvození; 20–32 blokové úlohy; 32–42 samostatně; 42–45 shrnutí.\par
\textbf{Návaznost:} V následující hodině budeme na každé pozici vybírat znak nezávisle: variace s opakováním.\par
\vspace{2mm}\noindent{\footnotesize\color{slate} \textbf{Východiska:} rozpis výuky MAT, 4. ročník, hodiny 29--41; tematický plán SPŠ; srovnání tematického rozsahu a náročnosti s realisticky.cz, oddíl Kombinatorika (lekce 090101--090119). Výklad, příklady a řešení jsou nově autorsky sestavené, nikoli převzatá sbírka.}\par
\end{document}